\documentclass[11pt,a4paper]{ctexart}
\usepackage[margin=2.2cm]{geometry}
\usepackage{amsmath,amssymb,amsthm,mathtools,booktabs,hyperref,microtype}
\hypersetup{colorlinks=true,linkcolor=black,urlcolor=blue}
\newtheorem{theorem}{定理}
\title{\bfseries 四次区间检测、有限层集合的稳定次数\\与剩余类核矩阵}
\author{研究札记 · F3D-STABLE}
\date{2026 年 9 月 29 日}
\begin{document}
\maketitle

\begin{abstract}
对
\[
F_{p,D}(n,d)=v_p(n+d)-v_p(n-d)
\]
的全域固定整数多项式实现，本文研究放大目标 \(kg\) 的最低次数增长率。首先在 \(p=3\) 给出任意长度有限整数区间的四次检测器，并证明增益 \(1,2\) 时四次最优。随后在射影剩余树上构造共同深度核矩阵 \(K_r\)，证明所有两端最终常值函数 \(g\) 的稳定次数
\[
\lim_{k\to\infty}\deg_p(kg)/k
\]
存在，并由一个有限矩阵逆的 \(\ell^1\) 范数精确给出。对有限层集合得到显式两两相互作用公式，连续区间在固定层数下恰好最省。状态为 \texttt{PAPER-AUDITED}，尚未 Lean 形式化。
\end{abstract}

\section{四次区间检测}
以下取 \(p=3\)。令
\[
u=n+d,\qquad v=n-d.
\]
固定
\[
L\ge1,\qquad q\in\{1,2\},
\]
定义
\[
A_{L,q}
=
(9u^2+3^qv^2)
(u^2+3^{2L-2}v^2),
\]
\[
B_{L,q}
=
(9u^2+v^2)
(u^2+3^{2L-2+q}v^2).
\]
输出
\[
\mathcal I_{L,q}
=
(A_{L,q}+B_{L,q},A_{L,q}-B_{L,q}).
\]

对
\[
x=u/v,\qquad t=v_3(x),
\]
置
\[
\theta_q(x)=\frac{x^2+3^q}{x^2+1}.
\]
当 \(q=1,2\) 时，
\[
v_3(\theta_q(x))
=
q\,\mathbf1_{\{t\ge1\}}.
\]
所以
\[
v_3\!\left(
\frac{\theta_q(3x)}
{\theta_q(3^{1-L}x)}
\right)
=
q(\mathbf1_{\{t\ge0\}}-\mathbf1_{\{t\ge L\}}).
\]
清分母后正是
\[
A_{L,q}/B_{L,q}.
\]
因此
\[
\boxed{
F_{3,D}(\mathcal I_{L,q}(n,d))
=
q\,\mathbf1_{\{0\le F_{3,D}(n,d)<L\}}.
}
\]

线性平移把它移动到任意有限整数区间
\[
[a,a+L-1].
\]

若次数至多 \(3\)，两端输出为零会迫使候选有理函数的根集中在单一赋值层，因此不能支持长度至少 \(2\) 的区间；单层情况再由模 \(3\) 两个单位类和 Hensel 简单根提升排除。故
\[
\boxed{
a_{3,I}(1)=a_{3,I}(2)=4.
}
\]

\section{射影剩余类核}
固定深度 \(r\)。射影直线
\[
\mathbf P^1(\mathbb Q_p)
\]
有
\[
N_r=(p+1)p^{r-1}
\]
个深度 \(r\) 剩余类。

定义共同深度核
\[
(K_r)_{ij}
=
\min\{v_p(u_iv_j-u_jv_i),r\}.
\]
令 \(P_j\) 是对深度 \(j\) 大类取平均的嵌套投影，则
\[
K_r
=
\sum_{j=1}^r
p^{r-j}P_j.
\]

对 \(\theta\ge0\)，设
\[
\lambda_j(\theta)
=
\frac{p^{r-j+1}-1}{p-1}+\theta.
\]
由嵌套投影谱分解，
\[
\boxed{
(K_r+\theta I)^{-1}
=
\lambda_r^{-1}I
-
\sum_{j=1}^{r-1}
(\lambda_{j+1}^{-1}-\lambda_j^{-1})P_j.
}
\]

特别取
\[
\tau=\frac1{p-1},\qquad
M_r=K_r+\tau I.
\]
则 \(M_r\) 正定，\(M_r^{-1}\) 非对角元非正，且
\[
M_r\mathbf1
=
\frac{p^r}{p-1}\mathbf1.
\]

\section{稳定次数公式}
设
\[
g:\mathbb Z\to\mathbb Z
\]
两端最终常值。取足够大的 \(r\)，让
\[
g(v_p(x))
\]
在每个深度 \(r\) 剩余类上恒定，形成向量 \(\mathbf g\)。记
\[
\bar g=\frac1{N_r}\sum_i g_i.
\]

则极限存在：
\[
\boxed{
\sigma_p(g)
=
\lim_{k\to\infty}
\frac{\deg_p(kg)}k.
}
\]
而且
\[
\boxed{
\sigma_p(g)
=
\frac12
\left\|
\left(K_r+\frac1{p-1}I\right)^{-1}
(\mathbf g-\bar g\mathbf1)
\right\|_1.
}
\]

下界的核心是：任一代数根对剩余类平均的射影弦距剖面都可写成
\[
M_r\mu_\alpha
\]
其中
\[
\mu_\alpha\ge0,\qquad
\|\mu_\alpha\|_1\le1.
\]
因此一个分子根或分母根最多贡献一单位正/负权重，直接得到 \(\ell^1\) 次数下界。

上界则对核矩阵每一列构造二次无抵消多项式因子，解有限有理线性方程并清分母，得到真实整数多项式乘积。继续细分剩余树，使构造代价趋于上述下界。次可加性
\[
a(k+\ell)\le a(k)+a(\ell)
\]
把子序列逼近提升成完整极限。

\section{有限层集合}
取
\[
g=\mathbf1_S,
\qquad S\subset\mathbb Z
\]
有限。

二值目标使 \(\ell^1\) 范数简化为一个二次型，最终得到
\[
\boxed{
\sigma_p(S)
=
\frac{2(p-1)^2}{p(p+1)}|S|
-
\frac{2(p-1)^3}{p(p+1)}
\sum_{\substack{i<j\\i,j\in S}}
p^{-(j-i)}.
}
\]

对每个有限增益 \(k\)，统一有
\[
\boxed{
a_{p,S}(k)
\ge
\lceil k\sigma_p(S)\rceil.
}
\]

当 \(p=3\)：
\[
\boxed{
\sigma_3(S)
=
\frac23|S|
-
\frac43
\sum_{i<j\in S}3^{-(j-i)}.
}
\]

因此两个检测层越近，稳定次数中的共享节省越大；相互作用按层距指数衰减。

\section{连续区间最省}
固定
\[
S=\{t_1<\cdots<t_L\}.
\]
因为
\[
t_j-t_i\ge j-i,
\]
有
\[
\sum_{i<j}p^{-(t_j-t_i)}
\le
\sum_{i<j}p^{-(j-i)}.
\]
交互项前系数为负，所以总代价最小时必须让层距尽可能小。

于是
\[
\boxed{
\sigma_p(S)
\ge
\frac{2(p-1)}{p+1}(1-p^{-L}),
}
\]
且 \(L\ge2\) 时等号当且仅当 \(S\) 是连续整数区间。

对区间
\[
[a,a+L-1]
\]
有显式闭式：
\[
\boxed{
\sigma_p([a,a+L-1])
=
\frac{2(p-1)}{p+1}(1-p^{-L}).
}
\]
特别地，
\[
\boxed{
\sigma_3([a,a+L-1])
=
1-3^{-L}.
}
\]

\section{阈值与边界}
单阈值
\[
\mathbf1_{\{t\ge1\}}
\]
的稳定代价是
\[
\boxed{
\frac{p-1}{p+1}.
}
\]
所以长有限区间的稳定代价趋于两个边界阈值代价之和。

稳定次数是增益 \(k\to\infty\) 后的次数/增益比；它不意味着每个有限 \(k\) 都达到取整下界，也不是运行时间或系数位长。

\begin{thebibliography}{9}
\bibitem{FiliPetsche} P. Fili and C. Petsche, work on energy integrals over local fields and projective chordal kernels.
\end{thebibliography}
\end{document}
